\chapter*{B}
\label{chap:b}

\term{Backward Stability}
A numerical algorithm is backward stable if its computed answer is the exact answer to a nearby problem. In symbols, the output for input $x$ is the exact result for $x+\Delta x$, where $\|\Delta x\|$ is small compared with $\|x\|$. \par\noindent\textbf{Applications:} backward stability explains why matrix algorithms can be trusted even when round-off error is unavoidable.

\term{Baire Category Theorem}
The Baire category theorem says that a complete metric space cannot be built as a countable union of thin sets, called nowhere-dense sets. Equivalently, a countable intersection of open dense sets is dense. \par\noindent\textbf{Applications:} it proves that certain objects exist in analysis and supports the uniform boundedness and open mapping theorems.

\term{Banach Fixed-Point Theorem}
If a map on a complete metric space brings points strictly closer together, then it has exactly one fixed point, and repeated iteration converges to it. This is also called the contraction mapping principle. \par\noindent\textbf{Applications:} it proves existence and uniqueness of solutions to equations and gives a practical iteration method for finding them.

\term{Banach Space}
A Banach space is a normed vector space in which every Cauchy sequence converges to a point in the space. Completeness means that a calculation whose successive approximations become arbitrarily close does not leave the space. \par\noindent\textbf{Applications:} Banach spaces provide the setting for infinite-dimensional approximation, integral equations, optimization, and differential equations.

\term{Banach--Steinhaus Theorem}
The uniform boundedness principle says that a family of bounded linear operators on a Banach space is uniformly bounded on the unit ball if each individual vector is sent to a bounded family of values. Its proof uses the Baire category theorem. \par\noindent\textbf{Applications:} it prevents hidden blow-up in families of approximation operators and is useful in functional analysis and numerical methods.

\term{Banach--Tarski Paradox}
Using a finite number of non-measurable pieces and rigid motions, a solid ball in three-dimensional space can be reassembled into two balls of the same size. This does not violate conservation of volume because the pieces have no well-defined volume; the construction uses the axiom of choice. \par\noindent\textbf{Applications:} the paradox clarifies the limits of assigning volume to arbitrary sets and the role of symmetry in measure theory.

\term{Base Case}
The base case is the first statement checked in a proof by mathematical induction. After it is established, an induction step shows that truth for one integer implies truth for the next. \par\noindent\textbf{Applications:} base cases are essential in recursive algorithms, proofs about sequences, and definitions of data structures.

\term{Basis}
A basis of a vector space is a linearly independent set that spans the whole space. Every vector has one and only one finite linear combination of basis vectors. \par\noindent\textbf{Applications:} coordinates, Fourier expansions, finite-element methods, and computer graphics all represent objects using a suitable basis.

\term{Bayes' Theorem}
Bayes' theorem relates conditional probabilities: $P(A\mid B)=P(B\mid A)P(A)/P(B)$ when $P(B)>0$. It updates the probability of a hypothesis after new evidence is observed. \par\noindent\textbf{Applications:} Bayesian statistics, medical diagnosis, spam filtering, fault detection, and machine learning use this update rule.

\term{Bernoulli Numbers}
The Bernoulli numbers $B_n$ are defined by $x/(e^x-1)=\sum_{n\ge0}B_nx^n/n!$. They occur in formulas for sums of powers and in the Euler--Maclaurin formula. \par\noindent\textbf{Applications:} they improve numerical integration and connect combinatorics with special values of the Riemann zeta function.

\term{Bernoulli Trial}
A Bernoulli trial is an experiment with two possible outcomes, usually called success and failure, with probabilities $p$ and $1-p$. Independent repetitions form a Bernoulli process. \par\noindent\textbf{Applications:} Bernoulli trials model yes/no tests, defective/nondefective items, clicks, infections, and binary classification outcomes.

\term{Berry's Paradox}
Berry's paradox arises from the phrase “the smallest positive integer not definable in fewer than twenty words.” The phrase appears to define such an integer in fewer than twenty words, creating a contradiction between informal language and formal definitions. \par\noindent\textbf{Applications:} it motivates precise languages in logic, computer science, and the study of what can be described or computed.

\term{Bertrand's Paradox}
Bertrand's paradox shows that the probability of a randomly chosen chord of a circle being longer than a side of an inscribed equilateral triangle depends on how “random chord” is defined. \par\noindent\textbf{Applications:} it teaches that a probability model must specify its sampling procedure, a principle important in statistics, simulation, and scientific experiments.

\term{Besov Space}
Besov spaces $B^s_{p,q}$ measure the smoothness of functions using how their detail changes across scales. The parameter $s$ describes smoothness, while $p$ and $q$ describe size and how different scales are combined; Sobolev spaces are important special cases. \par\noindent\textbf{Applications:} Besov spaces are used in partial differential equations, wavelet compression, image denoising, and nonlinear approximation.

\term{Bessel Function}
Bessel functions solve the differential equation $x^2y''+xy'+(x^2-n^2)y=0$. The functions $J_n$ and $Y_n$ are its standard independent solutions. \par\noindent\textbf{Applications:} they describe vibrations of circular membranes, heat flow in cylinders, electromagnetic waves, and acoustic modes.

\term{Bijection}
A bijection is a function that is both one-to-one and onto: every input has a different output, and every target is reached exactly once. A bijection proves that two sets have the same number of elements, even when they are infinite. \par\noindent\textbf{Applications:} bijections are used in counting, data encoding, coordinate changes, and proofs about infinite sets.

\term{Bilinear Form}
A bilinear form is a function $B:V\times W\to k$ that is linear in each input when the other is held fixed. Inner products and matrix expressions $x^TAy$ are examples. \par\noindent\textbf{Applications:} bilinear forms describe energy, angles, quadratic optimization, finite-element models, and the geometry of conic sections.

\term{Binary Operation}
A binary operation combines two elements of a set and returns an element of the same set: $*:S\times S\to S$. Addition and multiplication are examples, but the operation need not be commutative or associative. \par\noindent\textbf{Applications:} binary operations define the basic structure of groups, rings, logic gates, and algebraic data types.

\term{Binary Relation}
A binary relation between sets $A$ and $B$ is a collection of ordered pairs $(a,b)$ describing which elements are related. A function is a special relation in which each $a$ is paired with exactly one $b$. \par\noindent\textbf{Applications:} relations represent networks, database links, ordering, logical statements, and allowable transitions in algorithms.

\term{Binomial Coefficient}
The binomial coefficient $\binom nk=n!/[k!(n-k)!]$ counts the ways to choose $k$ objects from $n$ without regard to order. It is also the coefficient of $x^ky^{n-k}$ in $(x+y)^n$. \par\noindent\textbf{Applications:} binomial coefficients occur in probability, combinatorics, polynomial expansion, and error-correcting codes.

\term{Binomial Distribution}
The binomial distribution gives the probability of $k$ successes in $n$ independent Bernoulli trials: $P(X=k)=\binom nkp^k(1-p)^{n-k}$. Its parameters are the number of trials $n$ and success probability $p$. \par\noindent\textbf{Applications:} it models counts such as successful tests, conversions, defects, and disease cases in a fixed sample.

\term{Binomial Theorem}
The binomial theorem expands $(x+y)^n$ as $\sum_{k=0}^n\binom nkx^{n-k}y^k$. It converts a power of a sum into a sum of simpler terms. \par\noindent\textbf{Applications:} it is used in algebra, probability, numerical approximations, and the derivation of many combinatorial identities.

\term{Bipartite Graph}
A bipartite graph has its vertices divided into two sets, with every edge joining the two different sets. It has no odd-length cycle. \par\noindent\textbf{Applications:} bipartite graphs model job assignments, matching students to projects, network flow, recommendation systems, and molecular structures.

\term{Birch--Swinnerton-Dyer Conjecture}
For an elliptic curve $E$ over $\mathbb Q$, this conjecture relates the number of independent rational points on $E$ to how many times its $L$-function vanishes at $s=1$. It also predicts the leading coefficient using arithmetic quantities such as the regulator and Tate--Shafarevich group. \par\noindent\textbf{Applications:} it links analytic formulas with rational solutions and guides research in arithmetic geometry.

\term{Birthday Paradox}
In a group of 23 people, the probability that two share a birthday is already greater than $1/2$, assuming 365 equally likely birthdays. The probability of no match among $n$ people is $\prod_{k=0}^{n-1}(365-k)/365$. \par\noindent\textbf{Applications:} the effect is important in collision probabilities, hash tables, computer security, and the design of identifiers.

\term{Borel--Cantelli Lemmas}
If events $E_n$ satisfy $\sum P(E_n)<\infty$, then with probability one only finitely many of them occur. If the events are independent and $\sum P(E_n)=\infty$, then infinitely many occur with probability one. \par\noindent\textbf{Applications:} these lemmas test whether rare events eventually stop occurring and help prove convergence in probability theory.

\term{Borel Measure}
A Borel measure assigns sizes or probabilities to Borel sets, the sets generated from open sets by countable unions, intersections, and complements. On familiar spaces such as $\mathbb R^n$, it includes the usual notions of length, area, and volume. \par\noindent\textbf{Applications:} Borel measures form the foundation of probability, integration, and geometric measurement.

\term{Borel Set}
A Borel set is built from open sets using countable unions, countable intersections, and complements. The Borel sets form the smallest $\sigma$-algebra containing all open sets. \par\noindent\textbf{Applications:} intervals, regions, and countable combinations of them can be assigned probabilities or measures in analysis and statistics.

\term{Borsuk--Ulam Theorem}
Every continuous map from the sphere $S^n$ to $\mathbb R^n$ gives the same value at some pair of opposite points. For $n=2$, one consequence is that at some pair of antipodal points on Earth, temperature and atmospheric pressure are simultaneously equal. \par\noindent\textbf{Applications:} the theorem supports fair-division results such as the ham-sandwich theorem and studies of continuous data on spheres.

\term{Bound}
A bound is a number that limits a set, sequence, or function from above or below. A set is bounded if it has both an upper and a lower bound. \par\noindent\textbf{Applications:} bounds control error, guarantee stability, establish convergence, and quantify worst-case performance in algorithms.

\term{Boundary}
The boundary of a set $S$ consists of points whose every neighborhood meets both $S$ and its complement. For an interval $[0,1]$, the boundary is $\{0,1\}$. \par\noindent\textbf{Applications:} boundaries specify interfaces, constraints, and edge conditions in geometry, image processing, and differential equations.

\term{Boundary-value Problem}
A boundary-value problem asks for a solution of a differential equation subject to conditions on the boundary of a domain, such as fixed values (Dirichlet conditions) or fixed normal derivatives (Neumann conditions). \par\noindent\textbf{Applications:} it models temperature in a solid, vibration of a membrane, fluid flow, and electrostatic fields.

\term{Bounded Function}
A function is bounded if all its values lie between two fixed numbers: $m\le f(x)\le M$. For example, sine is bounded between $-1$ and $1$. \par\noindent\textbf{Applications:} boundedness prevents uncontrolled output and is used in stability analysis, probability, optimization, and signal processing.

\term{Bounded Linear Operator}
A linear map $T:X\to Y$ between normed spaces is bounded if $\|Tx\|\le C\|x\|$ for some constant $C$. For linear maps, boundedness is equivalent to continuity, and $\|T\|=\sup_{\|x\|\le1}\|Tx\|$ measures the largest amplification. \par\noindent\textbf{Applications:} operator norms quantify sensitivity in differential equations, numerical linear algebra, and signal transformations.

\term{Boy or Girl Paradox}
The probability that two children are both girls depends on how the information “at least one is a girl” was obtained. Under the usual equally likely interpretation it is $1/3$, while choosing a random child and finding that child is a girl gives $1/2$. \par\noindent\textbf{Applications:} the paradox teaches careful conditioning, which is essential in medical testing, surveys, and statistical reporting.

\term{Braess's Paradox}
Adding a road to a congested network can make every traveller's journey longer because selfish route choices can create a worse equilibrium. \par\noindent\textbf{Applications:} the paradox informs traffic planning, communication networks, supply chains, and the design of systems where local decisions affect global performance.

\term{Branch Point}
A branch point is a point around which a multi-valued complex function, such as $\sqrt z$ or $\log z$, cannot be made single-valued and continuous without choosing a branch cut. Going once around the point may change the value. \par\noindent\textbf{Applications:} branch points occur in complex integration, wave propagation, fluid flow, and mathematical physics.

\term{Brouwer Fixed-Point Theorem}
Every continuous map from a closed ball in $\mathbb R^n$ to itself has at least one fixed point $f(x)=x$. Continuity and the fact that the map stays inside the ball are essential. \par\noindent\textbf{Applications:} the theorem proves equilibrium existence in economics and game theory and helps establish solutions of nonlinear equations.

\term{Brownian Motion}
Brownian motion is a continuous random process whose independent increments are normally distributed and whose variance grows with elapsed time. It is the mathematical model of random diffusion. \par\noindent\textbf{Applications:} it models molecular motion, particle diffusion, financial price models, heat transport, and stochastic differential equations.

\term{Burali--Forti Paradox}
The collection of all ordinal numbers cannot be a set: if it were, it would have an ordinal larger than every ordinal in the collection, including itself. This contradiction helped motivate modern axioms that distinguish sets from proper classes. \par\noindent\textbf{Applications:} it clarifies the foundations of set theory and prevents uncontrolled definitions of “the set of everything.”

\term{Burnside's Lemma}
If a finite group $G$ acts on a finite set $X$, the number of distinct orbits is $|X/G|=|G|^{-1}\sum_{g\in G}|X^g|$, where $X^g$ is the set fixed by $g$. It counts objects while treating symmetric copies as the same. \par\noindent\textbf{Applications:} it counts necklaces, colorings, chemical structures, and configurations under rotation or reflection.

\term{Busy Beaver}
The Busy Beaver function gives the largest number of steps, or ones written, by any halting Turing machine with a fixed number of states. It grows faster than every computable function and therefore cannot itself be computed by an algorithm. \par\noindent\textbf{Applications:} Busy Beaver illustrates the limits of computation, undecidability, and the difference between simple programs and predictable behavior.
